% !PRINTED 208
\weberpage{208}
Ist $I$ eine \gesperrt{ursprüngliche} Function der $a$, d.\,h. eine solche,
deren Zahlenco\"efficienten keinen gemeinsamen Theiler haben, so können die
Co\"efficienten der Functionen $\Phi$ jedenfalls keinen anderen
gemeinschaftlichen Theiler haben, als eine Potenz von 3.

Was uns zu beweisen obliegt, ist, dass sie alle durch $3^\nu$ theilbar sind,
oder was dasselbe ist, dass, wenn wir $\Phi(A, D)$ als ursprüngliche Function
voraussetzen, $\nu = 0$ sein muss.

Es ist also die Frage: Können aus der ursprünglichen Function $\Phi(A, D)$
oder $B\Phi(A, D)$ Functionen mit dem Theiler 3 entstehen, wenn man die $A$,
$B$, $D$ durch ihre Ausdrücke in den $a$ ersetzt? Da $B$ ursprünglich ist, so
kann nach \S.\,2 die in den $a$ ausgedrückte Function $B\Phi(A, D)$ nur dann
den Theiler 3 haben, wenn ihn $\Phi(A, D)$ hat.

Also ist die Frage darauf zurückgeführt: Kann die ursprüngliche Function
$\Phi(A, D)$ den Theiler 3 erhalten, wenn $A$, $D$ durch ihre Ausdrücke
ersetzt werden?

Dass diese Frage verneint werden muss, können wir leicht so einsehen. Wir
denken uns zunächst in $\Phi(A, D)$ alle die Glieder beseitigt, deren
Co\"efficienten durch 3 theilbar sind; denn offenbar müssen auch die übrigen
Glieder noch dieselbe Eigenschaft behalten, durch Substitution der Ausdrücke
für $A$, $D$ den Theiler 3 zu erhalten.

Wir können zweitens annehmen, dass $\Phi(A, D)$ nicht den Factor $D$ hat,
denn durch Weglassen dieses Factors würde nach dem schon erwähnten Satze
(\S.\,2) die fragliche Eigenschaft nicht aufgehoben.

Dann aber müsste aus $\Phi(A, D)$ eine durch 3 theilbare Zahl entstehen, wenn
alle $a$, mit Ausnahme von $a_2$, gleich Null, $a_2 = 1$ gesetzt werden.
Dadurch wird $D = 0$, $A = 1$, und es müsste also $\Phi(1, 0)$ eine durch 3
theilbare Zahl sein; dies ist aber der Co\"efficient der höchsten Potenz von
$A$ in $\Phi(A, D)$, der nach Voraussetzung nicht durch 3 theilbar ist.
Damit haben wir also bewiesen:

\gesperrt{Jede ganzzahlige Invariante der biquadratischen Form lässt sich
rational und ganzzahlig durch $A$, $B$, $D$ darstellen, und zwar in einer der
beiden Formen}
\[
\Phi(A, D), \quad B\Phi(A, D).
\]
